
\subsubsection{3.1 Dataset}

Twenty pretrained convolutional neural networks were obtained from PyTorch torchvision. The models span three architecture families: ResNet (9 models), VGG (4 models), and DenseNet (4 models), plus AlexNet, SqueezeNet, and MobileNetV2. All models used ImageNet-1K pretraining with standard training recipes.

Models were classified into two depth categories: shallow (≤34 layers, n=10) and deep (≥50 layers, n=10). Shallow models included ResNet-18, ResNet-34, VGG-11, VGG-13, VGG-16, VGG-19, AlexNet, SqueezeNet1_0, MobileNetV2, and DenseNet-121. Deep models included ResNet-50, ResNet-101, ResNet-152, DenseNet-161, DenseNet-169, DenseNet-201, Wide-ResNet-50-2, Wide-ResNet-101-2, ResNeXt-50-32x4d, and ResNeXt-101-32x8d. The threshold was set to create balanced classes and avoid ambiguous intermediate depths (35-49 layers).

The dataset was split using stratified 80/20 division (16 training models, 4 test models) with fixed random seed 42. Test models were: alexnet (shallow), vgg13 (shallow), resnet152 (deep), and wide_resnet50_2 (deep).

\subsubsection{3.2 Feature Extraction}

Three feature types were extracted from model weights:

\textbf{Global Weight Statistics.} Layer-wise Frobenius norms were computed for all weight tensors, then aggregated into four statistics: mean, standard deviation, minimum, and maximum of layer-wise norms. For a weight tensor W, the Frobenius norm is ||W||_F = sqrt(sum(W_{ij}^2)).

\textbf{Layer-Wise Progression Features.} Six features were extracted: (1) linear regression slope of layer-wise norms versus layer position, (2) variance of layer-wise norms, (3) mean norm of first layer, (4) mean norm of last layer, (5) depth-weighted norm sum (sum of layer_position × norm), and (6) layer count.

\textbf{Architectural Features.} Eight structural properties were extracted: (1) residual block count (number of modules with downsample attribute), (2) dense connection count (layers containing "denselayer" in name), (3) bottleneck ratio (count of 1×1 convolutions divided by total convolutions), (4) total Conv2d layer count, (5) skip connection presence (binary indicator), (6) mean residual path norm (average norm of downsample layer weights), (7) transition layer count (DenseNet-specific), and (8) architecture family (one-hot encoded indicator).

\textbf{Batch Normalization Features.} Six features were extracted: (1) batch normalization layer count, (2) mean of gamma parameters, (3) standard deviation of gamma parameters, (4) mean of beta parameters, (5) standard deviation of beta parameters, and (6) depth-weighted batch normalization norm.

\subsubsection{3.3 Classification}

Logistic regression (scikit-learn implementation, C=1.0, solver='lbfgs', max_iter=1000, random_state=42) was used as the classifier. Features were normalized using StandardScaler (zero mean, unit variance) before training. The same classifier configuration was used across all feature types for controlled comparison.

\subsubsection{3.4 Validation Protocols}

\textbf{Random Initialization Test.} For each feature type, the same 20 model architectures were loaded with random initialization (pretrained=False in PyTorch), features were extracted, and a separate classifier was trained. This tests whether features reflect architectural properties (present in network structure regardless of training) versus training-induced properties (emergent from gradient flow). If random models achieve similar accuracy to pretrained models, features are architectural. If accuracy degrades substantially, features are training-induced.

\textbf{Within-Family Validation.} Separate classifiers were trained on ResNet-only (9 models) and DenseNet-only (4 models) subsets using stratified train-test splits. This controls for architecture family differences and tests whether depth signals exist within single families. VGG was excluded because all torchvision VGG variants are shallow (11-19 layers), providing no depth variation.

\subsubsection{3.5 Metrics}

Primary metric was test accuracy on the 4 held-out models. Secondary metrics included train accuracy (overfitting check), confusion matrix, and feature importance (logistic regression coefficients). For within-family validation, a threshold of ≥65\% accuracy was set as the success criterion. For overall classification, the initial target was ≥70\% test accuracy.
